% !TeX root = ../../Book3.tex
% 纲要标题：### 17.1 局部极小自由消解
% 纲要位置：计划纲要.md，第 2490 行。

\section{局部极小自由消解}
\label{b3:sec:1701}

自由消解中的一个单位系数把相邻两项的两个基向量配成一对，可以通过换基把这一对单独拆出。极小自由消解排除这样的配对，要求所有正次微分的矩阵元素都落在局部环的极大理想中。约化到剩余域后，微分同时变成零，各阶自由秩便能从模的同调中读取。

% 纲要内容（仅作写作计划，尚非教材正文）：
%
% * 极小生成元与极小微分
% * 单位系数消去及有限长度、各项有限秩消解的可缩直和分解
% * 剩余域张量后的简化
% * Betti 数及终止性
%

\subsection{极小生成元与极小微分}
\label{b3:subsec:1701:01}

给有限生成模写一个自由满射并不困难；困难在于分清哪些生成元确有必要，以及关系的关系会不会一直出现。局部环的剩余域把这两个问题变成向量空间问题。本节的自由消解各项均为有限秩自由模，但不预设消解长度有限，仍采用降次链复形记号。

\begin{definition}\label{b3:def:minimal-free-resolution}
设 \((R,\mathfrak m)\) 为交换 Noether 局部环，\(k=R/\mathfrak m\) 为其剩余域，\(M\) 为有限生成的 \(R\) 模。若增广正合复形
\[
\cdots\longrightarrow F_2\xrightarrow{\mathrm{d}_2}F_1
\xrightarrow{\mathrm{d}_1}F_0\xrightarrow{\epsilon}M\longrightarrow0
\]
中每个 \(F_i\) 为有限秩自由模，且 \(\mathrm{d}_i(F_i)\subseteq\mathfrak mF_{i-1}\) 对 \(i\ge1\) 成立，则称它为 \(M\) 的\term[jixiao-ziyou-xiaojie]{极小自由消解}[minimal free resolution]。
\end{definition}

正合性与 \(\mathrm{d}_1(F_1)\subseteq\mathfrak mF_0\) 说明增广映射诱导同构
\[
F_0/\mathfrak mF_0\xrightarrow{\sim}M/\mathfrak mM.
\]
因此极小性也包括零次项没有多余生成元。矩阵表示下，条件就是所有正次微分的矩阵元素都属于 \(\mathfrak m\)。

\ProseParagraphBreak
设 \(R\) 为交换环。若一个 \(R\) 模链复形的恒等链映射与零映射链同伦，则称它为\term[kesuo-fuxing]{可缩复形}[contractible complex]。相邻两次的恒等映射 \(R\xrightarrow{1}R\) 是最简单的例子；下面的命题说明它怎样从含单位系数的微分中出现。

\begin{proposition}[单位系数的消去]\label{b3:prop:unit-cancellation}
设 \((R,\mathfrak m)\) 为交换 Noether 局部环，\(k=R/\mathfrak m\)，\(M\) 为有限生成的 \(R\) 模，\(F_\bullet\to M\) 为各项均为有限秩自由模的自由消解。若某个整数 \(j\ge1\) 的微分 \(\mathrm{d}_j:F_j\to F_{j-1}\) 的矩阵含单位系数，则通过换基有增广复形的直和分解
\[
F_\bullet\cong G_\bullet\oplus\mathbb D(j),\qquad
\mathbb D(j):\quad0\longrightarrow R\xrightarrow{1}R\longrightarrow0,
\]
其中 \(\mathbb D(j)\) 的两个非零项位于次数 \(j,j-1\)，增广为零，它为可缩复形；\(G_\bullet\to M\) 仍为各项均为有限秩自由模的自由消解。由此，\(F_\bullet\) 极小等价于 \(F_\bullet\otimes_Rk\) 的全部正次微分为零，也等价于不能拆出任何 \(\mathbb D(j)\) 这样的非零直和项。
\end{proposition}
\begin{proof}
把单位系数交换到矩阵左上角，乘以其逆将它化为 \(1\)，再用可逆行、列变换清除同行、同列的其他系数，得到 \(\operatorname{diag}(1,D)\)。因此可选分解
\[
\begin{aligned}
F_j&=Ru\oplus F'_j,& F_{j-1}&=Rv\oplus F'_{j-1},\\
\mathrm{d}_j(u)&=v,& \mathrm{d}_j(F'_j)&\subseteq F'_{j-1}.
\end{aligned}
\]
由 \(\mathrm{d}_j\mathrm{d}_{j+1}=0\)，\(\mathrm{d}_{j+1}\) 的像在 \(u\) 方向的坐标必须为零。若 \(j>1\)，\(\mathrm{d}_{j-1}\mathrm{d}_j=0\) 给出 \(\mathrm{d}_{j-1}(v)=0\)；若 \(j=1\)，则由 \(\epsilon\mathrm{d}_1=0\) 得到 \(\epsilon(v)=0\)。所以 \(Ru\xrightarrow{1}Rv\) 与其余各项分离，给出所述直和分解。取 \(h_{j-1}(v)=u\)，其余同伦分量为零，便有 \(\mathrm{d}h+h\mathrm{d}=\operatorname{id}_{\mathbb D(j)}\)。此两项复形正合，因此余下的增广复形 \(G_\bullet\to M\) 也正合。

张量剩余域后微分为零，恰好表示原微分的全部矩阵系数属于 \(\mathfrak m\)。局部环中 \(\mathfrak m\) 之外的元素都是单位，所以不极小时可以作上述消去；极小时则在任何基下都不能出现单位系数，也不能有这样的两项直和因子。
\end{proof}

\begin{corollary}\label{b3:cor:finite-resolution-splitting}
设 \((R,\mathfrak m)\) 为交换 Noether 局部环，\(M\) 为有限生成的 \(R\) 模。给定 \(M\) 的有限长度、各项均为有限秩自由模的自由消解 \(F_\bullet\to M\)，可经有限次单位系数消去得到
\[
F_\bullet\cong F_\bullet^{\min}\oplus
\bigoplus_{\nu=1}^s\mathbb D(j_\nu),
\]
其中 \(F_\bullet^{\min}\to M\) 是极小自由消解，\(s\) 为非负整数，每个 \(\mathbb D(j_\nu)\) 是在次数 \(j_\nu,j_\nu-1\) 放置 \(R\xrightarrow{1}R\) 的两项复形。
\end{corollary}
\begin{proof}
原消解只有有限多个非零项，各项秩又有限，故总自由秩为有限整数。每次应用\cref{b3:prop:unit-cancellation}使总自由秩减少二，所以只需有限次消去。停止后所有正次微分均没有单位系数，即其系数全在 \(\mathfrak m\) 中，余下的消解正是极小自由消解。
\end{proof}

\begin{proposition}\label{b3:prop:minimal-resolution-existence}
设 \((R,\mathfrak m)\) 为交换 Noether 局部环，\(k=R/\mathfrak m\)，\(M\) 为有限生成的 \(R\) 模。则 \(M\) 有极小自由消解。它可通过依次选取满射 \(F_i\to K_i\) 构造，其中 \(K_0=M\)、\(K_{i+1}=\ker(F_i\to K_i)\)，且每个 \(F_i/\mathfrak mF_i\to K_i/\mathfrak mK_i\) 都是 \(k\) 线性同构。
\end{proposition}
\begin{proof}
取 \(M/\mathfrak mM\) 的一组 \(k\) 基并提升到 \(M\)。由\cref{b3:cor:local-generators}，提升构成极小生成组，给出 \(F_0\to M\)。因 \(R\) Noether 且 \(F_0\) 有限生成，核 \(K_1\) 也有限生成。对 \(K_1\) 重复这一过程，得到 \(F_1\to K_1\)，再合成到 \(F_0\)。

一般地令 \(K_0=M\)，依次取按 \(K_i/\mathfrak mK_i\) 的一组基提升所得的极小满射 \(\pi_i:F_i\to K_i\)，核为 \(K_{i+1}\)。所选基给出同构
\[
\overline{\pi_i}:F_i/\mathfrak mF_i\xrightarrow{\sim}K_i/\mathfrak mK_i.
\]
若 \(z\in K_{i+1}\)，则 \(\pi_i(z)=0\)，故 \(\overline{\pi_i}(z+\mathfrak mF_i)=0\)；由上式的单射性得 \(z\in\mathfrak mF_i\)。因此 \(K_{i+1}\subseteq\mathfrak mF_i\)，接起这些短正合列后，每个正次微分的像都属于下一项的 \(\mathfrak m\) 倍，得到所需极小消解。若某个核为零，以后全取零项。
\end{proof}

例如，设 \(k\) 为域，在 \(R=k[x,y]_{(x,y)}\) 上，模 \(R/(x,y)\) 有极小消解
\[
0\longrightarrow R\xrightarrow{\binom{-y}{x}}R^2
\xrightarrow{(x\quad y)}R\longrightarrow k\longrightarrow0.
\]
这也是已证的二元正则序列 Koszul 消解。在任何相邻的非负次数额外直加 \(R\xrightarrow{1}R\)，仍得到一个自由消解。这种加入可缩直和项的操作与\cref{b3:prop:unit-cancellation}中的消去互为逆过程。

\subsection{剩余域张量后的简化}
\label{b3:subsec:1701:02}

设 \(R\) 为交换环，\(M,N\) 为 \(R\) 模。本章所需的 Tor、Ext 来自自由消解。具体地，取 \(M\) 的自由消解 \(F_\bullet\)，约定
\[
\operatorname{Tor}_i^R(M,N)=H_i(F_\bullet\otimes_RN),\qquad
\operatorname{Ext}_R^i(M,N)=H^i\operatorname{Hom}_R(F_\bullet,N).
\]
由\cref{b3:lem:free-resolution-comparison}，消解之间的比较映射唯一到链同伦，故以上同调不依赖消解的选择。\cref{b3:lem:ext-basic-properties,b3:lem:tor-basic-properties}给出长正合列及 Tor 的平衡计算。对于极小消解，微分的像落在极大理想乘以相应自由模之中；取剩余域系数后，微分全部变为零，同调因而直接由消解各项得到。

\begin{proposition}\label{b3:prop:minimal-resolution-betti}
设 \((R,\mathfrak m)\) 为交换 Noether 局部环，\(k=R/\mathfrak m\)，\(M\) 为有限生成的 \(R\) 模，\(F_\bullet\to M\) 为其极小自由消解。则对每个整数 \(i\ge0\)，
\[
\operatorname{Tor}_i^R(M,k)\cong F_i\otimes_Rk,
\qquad
\operatorname{Ext}_R^i(M,k)\cong\operatorname{Hom}_R(F_i,k).
\]
任意两个极小自由消解同构为增广复形。
\end{proposition}
\begin{proof}
微分矩阵的元素在 \(\mathfrak m\) 中，张量 \(k\) 或映入 \(k\) 后全部成为零，所以同调就是相应项本身。
设 \(F_\bullet,F'_\bullet\) 为两个极小消解。由\cref{b3:lem:free-resolution-comparison}，取提升 \(M\) 的恒等映射的双向比较映射 \(f:F_\bullet\to F'_\bullet\)、\(g:F'_\bullet\to F_\bullet\)。两方向的复合各自与恒等链同伦；模掉 \(\mathfrak m\) 后，微分为零，同伦公式的右边也为零。因此
\[
\bar g_i\bar f_i=\operatorname{id}_{F_i/\mathfrak mF_i},
\qquad
\bar f_i\bar g_i=\operatorname{id}_{F'_i/\mathfrak mF'_i}.
\]
所以 \(\bar f_i:F_i/\mathfrak mF_i\to F'_i/\mathfrak mF'_i\) 是 \(k\) 线性同构，两边的自由秩相等。若共同秩为零，\(f_i\) 已是零模之间的同构；否则选基后 \(f_i\) 是方阵，且 \(\det(\bar f_i)\ne0\)，于是 \(\det(f_i)\notin\mathfrak m\)。局部环中极大理想之外的元素为单位，伴随矩阵公式因而给出 \(f_i\) 的逆。各次 \(f_i\) 可逆，且与微分和增广相容，构成增广复形同构。
\end{proof}

\subsection{Betti 数与消解终止性}
\label{b3:subsec:1701:03}

\begin{definition}\label{b3:def:local-betti-numbers}
设 \((R,\mathfrak m)\) 为交换 Noether 局部环，\(k=R/\mathfrak m\)，\(M\) 为有限生成的 \(R\) 模，\(i\ge0\) 为整数。整数
\(\beta_i^R(M)\coloneqq\dim_k\operatorname{Tor}_i^R(M,k)\)
称为 \(M\) 的第 \(i\) 个\term[betti-shu]{Betti 数}[Betti number]。
\end{definition}

\(\beta_0\) 是最少生成元数，\(\beta_1\) 是极小展示中的最少关系数。更高的数统计关系之间的关系。它们依赖底环：同一个剩余域 \(k\) 作为 \(k\) 模只有 \(\beta_0=1\)，作为 \(k[x,y]_{(x,y)}\) 模则有 \((\beta_0,\beta_1,\beta_2)=(1,2,1)\)。

\begin{corollary}\label{b3:cor:betti-gap-termination}
设 \((R,\mathfrak m)\) 为交换 Noether 局部环，\(k=R/\mathfrak m\)，\(M\) 为有限生成的 \(R\) 模，\(i\ge0\) 为整数。若 \(\beta_i^R(M)=0\)，则 \(M\) 的极小自由消解从第 \(i\) 项起全部为零。\(i=0\) 时，这一条件等价于 \(M=0\)。
\end{corollary}
\begin{proof}
\cref{b3:prop:minimal-resolution-betti}给出 \(F_i=0\)。若 \(i>0\)，极小满射 \(F_i\to K_i\) 使 \(K_i=0\)，后续所有核和极小自由项也为零。\(i=0\) 时满射 \(F_0\to M\) 直接给出 \(M=0\)。
\end{proof}

作为对照，设 \(k\) 为域，令 \(R=k[\epsilon]/(\epsilon^2)\)。其剩余域具有消解
\[
\cdots\xrightarrow{\epsilon}R\xrightarrow{\epsilon}R
\xrightarrow{\epsilon}R\longrightarrow k\longrightarrow0.
\]
乘 \(\epsilon\) 的核与像均为 \((\epsilon)\)，所以它正合且极小；每个 Betti 数均为一。有限生成并不保证消解终止。
